\documentclass[10pt,a4paper]{article}

\usepackage{../preamble-paper}
\addbibresource{../ref.bib}
\AtEveryBibitem{\clearfield{note}}
\numberwithin{equation}{section}

% 目录样式：... 点连接标题与页码
\usepackage{tocloft}
\renewcommand{\cftsecleader}{\cftdotfill{\cftdotsep}}
\renewcommand{\cftsubsecleader}{\cftdotfill{\cftdotsep}}
\renewcommand{\cftdotsep}{4.5}
\renewcommand{\cftsecnumwidth}{2.3em}
% 减少目录内行间距
\setlength{\cftbeforesecskip}{4pt plus 1pt}
\setlength{\cftbeforesubsecskip}{2pt plus 1pt}
\setlength{\cftbeforesubsubsecskip}{1pt plus 0.5pt}
\setlength{\cftbeforeparaskip}{0pt}

% Custom binary operators for transported ring structure
\newcommand{\sqplus}{\mathbin{\boxplus}}
\newcommand{\sqtimes}{\mathbin{\boxtimes}}
\newcommand{\sqminus}{\mathbin{\boxminus}}
\newcommand{\sqdiv}{\mathbin{\boxed{/}}}

\fancyhead[LE,RO]{\thepage}
\fancyhead[RE]{\textbf{草稿}}
\fancyhead[LO]{\textit{}}

% 覆盖 polyglossia 中文节编号，使用纯数字+点格式（无"第""节""章"）
\titleformat{\section}{\Large\rmfamily\bfseries}{\thesection.}{1em}{}
\titleformat{\subsection}{\large\rmfamily\bfseries}{\thesubsection}{1em}{}
\titleformat{\subsubsection}{\normalsize\rmfamily\bfseries}{\thesubsubsection}{1em}{}
% paragraph 内嵌加粗 + 悬挂缩进 + 间距控制
\newif\ifparablock
\AddToHook{para/begin}{\ifparablock\vspace{1.5ex plus 0.3ex}\parablockfalse\fi}
\titleformat{\paragraph}[runin]{\normalsize\rmfamily\bfseries}{\theparagraph}{1em}{\global\hangindent=2.5em\global\hangafter=1\global\parablocktrue\relax}
\titlespacing{\paragraph}{0pt}{1.5ex plus 0.3ex minus 0.2ex}{0.8em}

\begin{document}


\subsection{Banach代数}

复Banach空间$(A,\|\cdot\|)$若是复代数，
且范数满足$\| a b\| \leq \|a\| \cdot \|a\|$，
称\textbf{Banach代数}(Banach algebra)。
范数不等式保证了乘法的连续，
即当$a_n \to a,\ b_n \to b$时：
$\|a_n b_n - ab\| 
\leq \|a_nb_n - ab_n\| + \|a b_n - ab\|
\leq \|a_n - a\| \|b_n\| + \|a \|\|b_n - b\| \to 0$。
Banach空间$H$上的线性算子$\mathcal B(H)$是含幺Banach代数。

对含幺Banach代数的元 
$\forall a \in A$，
其谱
$\sigma(a) = \{ \lambda \in \mathbb{C} \mid \lambda 1_A - x \notin A^\times \}$
是 $\mathbb C$ 的非空紧子集。


设若不然，假设 $\sigma(x) = \emptyset$。
这意味着对于所有的 $\lambda \in \mathbb{C}$，
特征算子 $\lambda 1_A - x$ 在 $A$ 中都是可逆的。
由此，我们可以定义全局有定义的预解算子值函数：
$R(\lambda) = (\lambda 1_A - x)^{-1}, \quad \forall \lambda \in \mathbb{C}$

任取一个有界线性泛函 $\varphi \in A^*$，并构造一个复值标量函数 $f: \mathbb{C} \to \mathbb{C}$：
\begin{equation}
f(\lambda) = \varphi(R(\lambda)) = \varphi\left( (\lambda 1_A - x)^{-1} \right)
\end{equation}

由于巴拿赫代数中的求逆映射是解析的，
函数 $f(\lambda)$ 在整个复平面 $\mathbb{C}$ 上处处全纯（即 $f$ 是一个整函数）。

接下来我们分析 $R(\lambda)$ 在 $|\lambda| \to \infty$ 时的渐近行为：
\begin{equation}
\|R(\lambda)\| = \left\| \lambda^{-1} \left(1_A - \frac{x}{\lambda}\right)^{-1} \right\| \le \frac{1}{|\lambda|} \frac{1}{1 - \frac{\|x\|}{|\lambda|}} \to 0 \quad (\text{当 } |\lambda| \to \infty)
\end{equation}
因此，$\lim_{|\lambda| \to \infty} f(\lambda) = 0$。根据连续性，整函数 $f$ 在整个复平面上必定是有界的。

根据复分析中的 \textbf{Liouville 定理}，复平面上任何有界整函数都必定是常数。因此，$f(\lambda)$ 是常数函数。又因为其无穷远极限为零，迫使对所有的 $\lambda \in \mathbb{C}$ 都有 $f(\lambda) = 0$。

由于这一结论对任意有界线性泛函 $\varphi \in A^*$ 都成立，根据 Hahn-Banach 定理（对偶泛函的非退化性），我们不得不得出：
\begin{equation}
R(\lambda) = (\lambda 1_A - x)^{-1} = 0_A, \quad \forall \lambda \in \mathbb{C}
\end{equation}
这产生了一个不可调和的矛盾，因为一个可逆元素绝不可能是零元素。因此，谱 $\sigma(x)$ 必定非空。
\end{proof}

\begin{proof}[定理~\ref{thm:gelfand-mazur} 的证明]
任取元素 $x \in A$。根据引理~\ref{lem:nonempty-spectrum}，其谱 $\sigma(x)$ 非空。因此，至少存在一个复数 $\lambda_x \in \sigma(x)$。

根据谱的定义，对应的特征算子
\begin{equation}
\lambda_x 1_A - x
\end{equation}
在 $A$ 中是不可逆的。

由于 $A$ 被假设为一个除代数（非零元皆可逆），除代数中唯一不具备逆元的元素只能是零元素 $0_A$。因此，我们被迫得出：
\begin{equation}
\lambda_x 1_A - x = 0_A \implies x = \lambda_x 1_A
\end{equation}
这说明 $A$ 中的每一个元素 $x$，本质上都只是单位元 $1_A$ 的一个复标量倍数。

我们现在定义如下规范映射：
\begin{equation}
\Phi: \mathbb{C} \to A, \quad \Phi(\lambda) = \lambda 1_A
\end{equation}
\begin{enumerate}[label=(\roman*)]
  \item \textbf{代数同态}：显然，$\Phi$ 保持了加法、标量乘法和乘法（因为 $\Phi(\alpha\beta) = \alpha\beta 1_A = (\alpha 1_A)(\beta 1_A) = \Phi(\alpha)\Phi(\beta)$）。
  \item \textbf{双射性}：由于 $1_A \neq 0_A$，$\lambda 1_A = 0_A \implies \lambda = 0$，故 $\Phi$ 是单射。又由前面的推导，任何 $x \in A$ 都可以写成 $\lambda_x 1_A = \Phi(\lambda_x)$ 的形式，故 $\Phi$ 也是满射。
  \item \textbf{等距性}：对任意 $\lambda \in \mathbb{C}$，其范数满足：
  \begin{equation}
  \|\Phi(\lambda)\| = \|\lambda 1_A\| = |\lambda| \|1_A\| = |\lambda|
  \end{equation}
  因为含幺赋范代数的单位元满足规范化条件 $\|1_A\| = 1$。
\end{enumerate}
综上所述，$\Phi$ 是复巴拿赫代数之间的一个等距同构映射。
\end{proof}

\begin{remark}[几何重构的本质]
Gelfand-Mazur 定理解释了为什么复数域 $\mathbb{C}$ 在经典和量子代数重构中具有统治地位。

设 $A$ 是一个交换复巴拿赫代数，$\mathfrak{m} \subset A$ 是它的一个极大理想。由于 $\mathfrak{m}$ 是闭的，商代数 $A/\mathfrak{m}$ 继承了巴拿赫代数的结构。并且，由于 $\mathfrak{m}$ 的极大性，商环 $A/\mathfrak{m}$ 是一个除代数。

根据 Gelfand-Mazur 定理，商巴拿赫除代数必定与复数域等距同构：
\begin{equation}
A/\mathfrak{m} \cong \mathbb{C}
\end{equation}
这一同构极具物理和几何威能：它确保了每一个极大理想 $\mathfrak{m}$，都唯一地对应着一个乘性线性泛函（即特征标） $\phi \in \Omega(A)$ 使得 $\ker(\phi) = \mathfrak{m}$。这正是“点”能从“纯代数”中重新被完好无损地提炼出来的最高代数保证。
\end{remark}



\end{document}